% SOURCE: written for the TMLR submission. Kept short and specific: the two
% concerns that are actually particular to THIS paper are naming real
% repositories and describing an exploitable procedure, so those are what the
% statement addresses rather than generic ML-risk boilerplate.
\section*{Broader Impact Statement}

This paper studies how experiments are measured rather than proposing a model,
a dataset, or a deployed system, so its effect is on research practice. Two
consequences deserve naming. First, the audit names repositories and commits:
we report structural preconditions and, where we re-executed, realized
divergence, but claim no published result is wrong. These are careful
implementations, written when nothing in the tooling or the review process
asked for probe-RNG isolation, and the same reading returned the same verdict
on our own predecessor first. Second, a procedure that inflates a reported
gap is a recipe as well as a diagnosis. Disclosure is safer: the mechanism is
already live in widely used released code, so it operates whether or not
practitioners intend it. Isolation is often a small code change, but the full
repair also requires deterministic execution and per-path repeat checks,
with runtime, storage and platform costs.
